% Generated by build_field_comparison_atlas.jl.
\begingroup
\setlength{\fboxsep}{1.0mm}
\setlength{\fboxrule}{.4pt}
\newcommand{\FieldRayCard}[7]{%
  \fcolorbox{#6}{white}{%
    \begin{minipage}[t][6.10cm][t]{9.40cm}
      \vspace{0pt}%
      \colorbox{#6!12}{\parbox{\dimexpr\linewidth-2\fboxsep\relax}{%
        \fontsize{12}{13}\selectfont\bfseries $\mathsf{R}_{#1}$%
        \hfill\normalfont\fontsize{10}{11}\selectfont #2\quad #5}}
      \par\vspace{.6mm}
      {\centering\fontsize{9.5}{11}\selectfont
        $\mathbb{Q}:\ #3$\par $\mathbb{F}_{2}:\ #4$\par}
      \begin{minipage}[c][4.82cm][c]{\linewidth}\centering
        \resizebox{4.70cm}{!}{\input{#7}}%
      \end{minipage}
    \end{minipage}%
  }%
}
\newcommand{\FieldKeyCard}[2]{%
  \fcolorbox{cellline!50}{paperblue!22}{%
    \begin{minipage}[t][4.95cm][t]{12.80cm}
      \vspace{0pt}\fontsize{12}{14}\selectfont\bfseries #1\par\vspace{2mm}
      \normalfont\fontsize{10.5}{13}\selectfont #2
    \end{minipage}%
  }%
}
{\fontsize{24}{27}\selectfont\bfseries The 12 ray classes}\par\vspace{1.5mm}
{\fontsize{11}{13}\selectfont\color{muted}
  Minimal metric-tree representatives under the Weyl group of order 12.
  The 11 QQ classes persist; $\mathsf{R}_{12}$ is specific to GF(2).}\par\vspace{2.5mm}
\FieldRayCard{1}{12 rays}{\underline{1},\,4,\,23,\,24,\,38,\,39,\,42,\,43,\,48,\,49,\,50,\,51}{6,\,8,\,16,\,17,\,\underline{32},\,33,\,34,\,35,\,37,\,38,\,41,\,47}{Both fields}{cellline}{field_comparison_assets/rays/ray_01.tex}\hfill
\FieldRayCard{2}{6 rays}{\underline{2},\,25,\,26,\,28,\,31,\,32}{5,\,15,\,20,\,23,\,\underline{42},\,48}{Both fields}{cellline}{field_comparison_assets/rays/ray_02.tex}\hfill
\FieldRayCard{3}{6 rays}{\underline{3},\,5,\,16,\,27,\,46,\,47}{7,\,\underline{28},\,30,\,31,\,43,\,49}{Both fields}{cellline}{field_comparison_assets/rays/ray_03.tex}\hfill
\FieldRayCard{4}{6 rays}{\underline{6},\,7,\,11,\,12,\,40,\,41}{13,\,\underline{14},\,18,\,21,\,24,\,25}{Both fields}{cellline}{field_comparison_assets/rays/ray_04.tex}\par\vspace{1.5mm}
\FieldRayCard{5}{3 rays}{\underline{8},\,33,\,34}{1,\,\underline{19},\,22}{Both fields}{cellline}{field_comparison_assets/rays/ray_05.tex}\hfill
\FieldRayCard{6}{3 rays}{\underline{9},\,10,\,30}{\underline{2},\,26,\,27}{Both fields}{cellline}{field_comparison_assets/rays/ray_06.tex}\hfill
\FieldRayCard{7}{6 rays}{\underline{13},\,17,\,19,\,20,\,44,\,45}{10,\,11,\,45,\,51,\,53,\,\underline{54}}{Both fields}{cellline}{field_comparison_assets/rays/ray_07.tex}\hfill
\FieldRayCard{8}{3 rays}{\underline{14},\,21,\,22}{4,\,\underline{46},\,52}{Both fields}{cellline}{field_comparison_assets/rays/ray_08.tex}\par\vspace{1.5mm}
\FieldRayCard{9}{2 rays}{\underline{15},\,29}{29,\,\underline{40}}{Both fields}{cellline}{field_comparison_assets/rays/ray_09.tex}\hfill
\FieldRayCard{10}{3 rays}{\underline{18},\,35,\,36}{9,\,44,\,\underline{50}}{Both fields}{cellline}{field_comparison_assets/rays/ray_10.tex}\hfill
\FieldRayCard{11}{1 ray}{\underline{37}}{\underline{12}}{Both fields}{cellline}{field_comparison_assets/rays/ray_11.tex}\hfill
\FieldRayCard{12}{3 rays}{\text{absent}}{\underline{3},\,36,\,39}{\textbf{GF(2) only}}{rootrust}{field_comparison_assets/rays/ray_12.tex}\par\vspace{1.5mm}
\FieldKeyCard{Reading the trees}{%
  \begin{tikzpicture}[x=1cm,y=1cm]
    \AtlasTaxonOne[minimum size=6mm]{0,0}{1}
    \AtlasTaxonTwo[minimum size=8.6mm]{1.05,0}{2}
    \AtlasTaxonThree[minimum size=6mm]{2.1,0}{3}
    \AtlasTaxonFour[minimum size=6mm]{3.15,0}{4}
    \AtlasTaxonFive[minimum size=8.6mm]{4.2,0}{5}
    \AtlasTaxonSix[minimum size=6mm]{5.25,0}{6}
    \AtlasTaxonSeven{6.4,0}{.72}{7}
  \end{tikzpicture}\par\vspace{2mm}
  Taxa 1--3 are blue; 4--6 are gold; 7 is the red heart.
  Shapes distinguish the three paired labels.\par\vspace{2mm}
  Dark blue edges are bounded; gray edges are pendant.
  A, B, \ldots\ name internal vertices. Edge numbers are exact lengths;
  drawn lengths are schematic.\par\vspace{2mm}
  Ray trees use the primitive integral Coxeter-balanced metric with minimum
  pendant length zero. Cell trees are the sums of their normalized rays, normalized in the same way.
}\hfill
\FieldKeyCard{Fields and fan sizes}{%
  \begin{tabular}{@{}lrr@{}}
    & $\mathbb{Q}$ & $\mathbb{F}_{2}$\\[2mm]
    Rays & 51 & 54\\
    Ray orbits & 11 & 12\\[1mm]
    Maximal cells & 248 & 254\\
    Maximal-cell orbits & 31 & 32
  \end{tabular}\par\vspace{2mm}
  Both fans have dimension 6 and lineality dimension 3.\par\vspace{2mm}
  Lists: field-specific census ray indices. Underlining: pictured tree.
  Classes 1--11: QQ ray-orbit order.
}\hfill
\FieldKeyCard{The central ray $\mathsf{R}_{12}$}{%
  Its three GF(2) rays, $3,36,39$, are the centers of the three
  Dressian tetrahedra. No QQ ray belongs to this class.\par\vspace{2mm}
  In each tetrahedron, the four QQ boundary triangles are replaced by the
  six GF(2) triangles joining the center to the six edges.\par\vspace{2mm}
  Under the full Weyl group, these give three QQ-only maximal-cell orbits
  and four GF(2)-only maximal-cell orbits. The detailed colored comparison
  is on page 3.
}\par
\endgroup
